\chapterimage{chapter_pointers.jpg} % Chapter heading image

\chapter{Estruturas de Controle, Operadores Lógicos e Relacionais}

Até agora seus programas executam de forma linear: linha por linha, sem desvios. Mas e se você quisesse fazer algo diferente dependendo de uma condição? \say{Se o usuário digitar uma senha correta, deixa entrar. Senão, nega acesso.} Para isso existem as estruturas de controle.

\section{Operadores relacionais}

Operadores relacionais comparam dois valores e retornam verdadeiro (true) ou falso (false). São eles:

\begin{itemize}
  \item \textbf{==} (igual): verifica se dois valores são iguais
  \item \textbf{!=} (diferente): verifica se dois valores são diferentes
  \item \textbf{<} (menor): verifica se o primeiro é menor que o segundo
  \item \textbf{>} (maior): verifica se o primeiro é maior que o segundo
  \item \textbf{<=} (menor ou igual): verifica se o primeiro é menor ou igual ao segundo
  \item \textbf{>=} (maior ou igual): verifica se o primeiro é maior ou igual ao segundo
\end{itemize}

Exemplos:

\begin{ccode}
  int idade = 18;
  int permitido = 21;

  idade == permitido      // false (18 não é igual a 21)
  idade != permitido      // true (18 é diferente de 21)
  idade < permitido       // true (18 é menor que 21)
  idade > permitido       // false (18 não é maior que 21)
  idade <= permitido      // true (18 é menor ou igual a 21)
  idade >= permitido      // false (18 não é maior ou igual a 21)
\end{ccode}

\begin{quote}
\textit{Cuidado com ==! Sério, esse é o erro mais comum desde a invenção do pão de queijo.}
\end{quote}

Um erro comum é usar uma única igualdade (\textit{=}) em vez de duas (\textit{==}). É tipo confundir um casamento com um namorado. A primeira é atribuição (você está \textit{dando} algo), a segunda é comparação (você está \textit{perguntando} se são iguais).

\begin{ccode}
  if (x = 5) { }    // ERRADO: atribui 5 a x (e seu código fica estranho)
  if (x == 5) { }   // CORRETO: compara x com 5 (e todos ficarão felizes)
\end{ccode}

Esse erro é tão comum que compiladores têm warnings especiais para ele. Se você cometer esse erro, o compilador vai olhar para você tipo aquele professor que te pega copiando a lição de casa.

\section{Operadores lógicos}

Operadores lógicos combinam múltiplas condições:

\subsection{E (&&)}

Retorna verdadeiro apenas se AMBAS as condições forem verdadeiras.

\begin{ccode}
  int idade = 25;
  int tem_carteira = 1;

  (idade >= 18) && (tem_carteira == 1)  // true (ambas condições verdadeiras)
\end{ccode}

\subsection{OU (||)}

Retorna verdadeiro se PELO MENOS UMA das condições for verdadeira.

\begin{ccode}
  int eh_policial = 1;
  int eh_seguranca = 0;

  (eh_policial == 1) || (eh_seguranca == 1)  // true (primeira é verdadeira)
\end{ccode}

\subsection{NÃO (!)}

Inverte o resultado de uma condição. Verdadeiro vira falso e vice-versa.

\begin{ccode}
  int chovendo = 1;
  !chovendo  // false (inverte 1)

  int ensolarado = 0;
  !ensolarado  // true (inverte 0)
\end{ccode}

\section{O comando if}

O comando \textit{if} (se) executa um bloco de código apenas se uma condição for verdadeira.

Sintaxe:

\begin{ccode}
  if (condição) {
    // código executado se condição for verdadeira
  }
\end{ccode}

Exemplo:

\begin{ccode}
  int nota = 70;

  if (nota >= 60) {
    printf("Aluno aprovado!\n");
  }
\end{ccode}

Se a nota fosse menor que 60, a mensagem não seria impressa.

\subsection{O comando else}

O comando \textit{else} (senão) executa um bloco alternativo se a condição for falsa.

\begin{ccode}
  int nota = 50;

  if (nota >= 60) {
    printf("Aluno aprovado!\n");
  } else {
    printf("Aluno reprovado!\n");
  }
\end{ccode}

Neste caso, como nota é 50, a segunda mensagem será impressa.

\subsection{O comando else if}

Quando você tem múltiplas condições, use \textit{else if} (senão se):

\begin{ccode}
  int nota = 75;

  if (nota >= 90) {
    printf("Excelente!\n");
  } else if (nota >= 80) {
    printf("Muito bom!\n");
  } else if (nota >= 70) {
    printf("Bom!\n");
  } else if (nota >= 60) {
    printf("Satisfatório!\n");
  } else {
    printf("Insuficiente!\n");
  }
\end{ccode}

As condições são avaliadas em ordem, de cima para baixo. Assim que uma for verdadeira, seu bloco é executado e o resto é ignorado.

\subsection{Ifs aninhados}

Você pode colocar \textit{ifs} dentro de \textit{ifs}:

\begin{ccode}
  int idade = 25;
  int tem_carteira = 1;

  if (idade >= 18) {
    if (tem_carteira == 1) {
      printf("Pode dirigir!\n");
    } else {
      printf("Deve tirar carteira!\n");
    }
  } else {
    printf("Muito jovem para dirigir!\n");
  }
\end{ccode}

Porém, ifs aninhados podem ficar confusos. Frequentemente é melhor usar operadores lógicos (\textit{\&\&} ou \textit{||}).

\section{O comando switch}

O \textit{switch} é útil quando você tem uma variável que pode ter vários valores distintos.

Sintaxe:

\begin{ccode}
  switch (variavel) {
    case valor1:
      // código se variavel == valor1
      break;
    case valor2:
      // código se variavel == valor2
      break;
    default:
      // código se nenhum case corresponder
      break;
  }
\end{ccode}

Exemplo prático:

\begin{ccode}
  int dia_semana = 3;

  switch (dia_semana) {
    case 1:
      printf("Segunda\n");
      break;
    case 2:
      printf("Terça\n");
      break;
    case 3:
      printf("Quarta\n");
      break;
    case 4:
      printf("Quinta\n");
      break;
    case 5:
      printf("Sexta\n");
      break;
    case 6:
      printf("Sábado\n");
      break;
    case 7:
      printf("Domingo\n");
      break;
    default:
      printf("Dia inválido\n");
      break;
  }
\end{ccode}

Este código imprime \say{Quarta} porque \textit{dia\_semana} é 3.

\begin{quote}
\textit{Nunca esqueça do break! Sério. Esse é o tipo de erro que faz você passar 3 horas debugando.}
\end{quote}

O \textit{break} interrompe a execução do \textit{switch}. É tipo sair de uma loja: você entra em um case, faz suas compras (executa o código), e depois sai. Sem \textit{break}, você é aquele cliente estranho que segue andando pela loja mesmo tendo conseguido o que queria. A execução continua para o próximo \textit{case}, mesmo que a condição não seja atendida (em outras palavras: caos):

\begin{ccode}
  // SEM break:
  switch (3) {
    case 3:
      printf("Quarta\n");
      // Falta break aqui!
    case 4:
      printf("Quinta\n");
      break;
  }
  // Imprime: Quarta
  //          Quinta
\end{ccode}

\section{Operador ternário (?:)}

O operador ternário é uma forma compacta de \textit{if-else} em uma única linha.

Sintaxe:

\begin{ccode}
  (condição) ? valor_se_verdadeiro : valor_se_falso
\end{ccode}

Exemplo:

\begin{ccode}
  int idade = 25;
  char resultado[10];

  if (idade >= 18) {
    strcpy(resultado, "adulto");
  } else {
    strcpy(resultado, "menor");
  }
\end{ccode}

Pode ser reescrito com operador ternário:

\begin{ccode}
  int idade = 25;
  char *status = (idade >= 18) ? "adulto" : "menor";
  printf("%s\n", status);  // Imprime: adulto
\end{ccode}

Use o operador ternário para atribuições simples. Para lógica mais complexa, continue usando \textit{if-else}.

Agora você domina as estruturas de controle em C. Com \textit{if}, \textit{else if}, \textit{switch} e operadores lógicos, você pode fazer seu programa tomar decisões inteligentes. No próximo capítulo, vamos aprender sobre repetições com loops.

%------------------------------------------------



